We report seven complementary metrics. Mean absolute error (MAE) and root mean squared error (RMSE) measure how close the point prediction is to the observed count at each cell-day, with RMSE penalizing large errors more heavily. The Wasserstein distance instead compares the empirical distribution of the observed counts with that of the predictions within each test window, so it rewards reproducing the overall mix of zero, low and high counts even when individual cell-days are misplaced. Because vessel activity is sparse, we also assess the detection of active cell-days ($y>0$) through accuracy and F1, the harmonic mean of precision and recall, which is more informative than accuracy under strong class imbalance. Finally, the mean log-likelihood of the observed counts and the continuous ranked probability score (CRPS), both proper scoring rules~\citep{gneiting2007strictly}, evaluate the whole predictive distribution rather than a point prediction: a Poisson distribution with the predicted rate for the baselines, and the hurdle distribution for Red-LGCP. The log-likelihood strongly penalizes assigning low probability to the observed count, whereas the CRPS, expressed in counts, generalizes the absolute error to probabilistic forecasts. Full formulas for all seven metrics are given in \S\ref{app:metrics}.
